% Luc Destombe --- licence CC BY-NC-SA 4.0
% https://creativecommons.org/licenses/by-nc-sa/4.0/deed.fr
% Fichier autonome : compiler avec pdflatex, deux fois (signets, nombre de pages).
\def\lyceeversion{1.1}
\documentclass[a4paper,10pt]{article}
\makeatletter
\newif\iflycee@niveaux
\newif\iflycee@palatino
\lycee@palatinotrue

\RequirePackage[utf8]{inputenc}
\RequirePackage[T1]{fontenc}
\RequirePackage[french]{babel}
\RequirePackage{etoolbox}

\newcommand{\ShorthandsOff}{\@ifundefined{shorthandoff}{}{\shorthandoff{;:!?}}}
\newcommand{\ShorthandsOn}{\@ifundefined{shorthandon}{}{\shorthandon{;:!?}}}
\ShorthandsOff
\AtBeginDocument{\ShorthandsOn}
\AtBeginEnvironment{tikzpicture}{\ShorthandsOff}

\RequirePackage{geometry}
\RequirePackage{fancyhdr}
\RequirePackage{lastpage}
\RequirePackage{multicol}
\RequirePackage{array}
\RequirePackage{enumitem}
\RequirePackage{xcolor}
\RequirePackage{needspace}

\RequirePackage{amsmath,amssymb}
\RequirePackage{mathpazo}
\renewcommand{\ttdefault}{lmtt}
\DeclareMathSizes{10.95}{10.95}{8}{5.5}
\begingroup\catcode`\_=\active \catcode`\^=\active
\gdef_#1{\sb{\scriptscriptstyle #1}}
\gdef^#1{\sp{\scriptscriptstyle\lycee@moinsepais #1}}
\endgroup
\newcommand\lycee@moins{\mathbin{%
  \setbox\z@\hbox{$\m@th\scriptscriptstyle\lycee@moinsorig$}%
  \dimen@=.55\wd\z@ \advance\dimen@-.5pt
  \hbox to.75\wd\z@{\hss$\m@th\scriptscriptstyle\vcenter{\hbox{%
    \tikz\draw[line width=.5pt,line cap=round](0,0)--(\the\dimen@,0);}}$\hss}}}
\begingroup\lccode`\~=`\- \lowercase{\endgroup\let~}\lycee@moins
\newcommand\lycee@moinsepais{\mathcode`\-="8000 }
\AtBeginDocument{\mathchardef\lycee@moinsorig=\mathcode`\-\relax}
\AtBeginDocument{\catcode`\_=12 \mathcode`\_="8000
  \catcode`\^=12 \mathcode`\^="8000 }
\newcommand\lycee@exposants{%
  \fontdimen13\textfont2=.48\fontdimen6\textfont2
  \fontdimen14\textfont2=.48\fontdimen6\textfont2
  \fontdimen15\textfont2=.48\fontdimen6\textfont2
  \fontdimen13\scriptfont2=.48\fontdimen6\scriptfont2
  \fontdimen14\scriptfont2=.48\fontdimen6\scriptfont2
  \fontdimen15\scriptfont2=.48\fontdimen6\scriptfont2}
\every@math@size\expandafter{\the\every@math@size\lycee@exposants}
\newlength{\retraitcalculs}
\setlength{\retraitcalculs}{2em}

\RequirePackage{tikz}
\usetikzlibrary{arrows.meta,calc,babel,shapes.misc}
\RequirePackage[most]{tcolorbox}
\tcbuselibrary{skins,breakable}

\definecolor{coulDef}{HTML}{1F4E79}
\definecolor{coulProp}{HTML}{7030A0}
\definecolor{coulRem}{HTML}{C55A11}
\definecolor{coulEx}{HTML}{2E7D32}
\definecolor{coulExo}{HTML}{B03A2E}
\definecolor{coulPart}{HTML}{1F4E79}
\definecolor{coulMeth}{HTML}{1B6E4F}
\definecolor{coulCourbe}{HTML}{0B5ED7}
\definecolor{coulCourbeB}{HTML}{C00000}
\definecolor{coulCourbeC}{HTML}{2E7D32}
\definecolor{coulGrille}{HTML}{8C8C8C}
\definecolor{coulRep}{HTML}{1B6E4F}
\definecolor{coulBar}{HTML}{2E7D32}

\newcommand{\R}{\mathbb{R}}
\newcommand{\Rs}{\mathbb{R}^{*}}
\renewcommand{\le}{\leqslant}
\renewcommand{\ge}{\geqslant}
\renewcommand{\approx}{\simeq}
\newcommand{\vect}[1]{\overrightarrow{#1}}
\newcommand{\norme}[1]{\left\lVert #1 \right\rVert}
\newcommand{\intervalle}[2]{\left[#1\,;#2\right]}
\RequirePackage{cancel}
\renewcommand{\CancelColor}{\color{coulExo}}
\newcommand{\fois}[1]{\times{\color{coulExo}#1}}
\newcommand{\barre}[1]{\cancel{#1}}

\tikzset{grille/.style={coulGrille,line width=0.4pt}}
\newcommand{\quadrillage}[4]{\draw[grille,step=1] (#1,#2) grid (#3,#4);}
\tikzset{
  croix/.style={cross out,draw,line width=0.6pt,inner sep=0pt,minimum size=4pt},
  croixdroite/.style={croix,rotate=45,line width=0.8pt},
}
\newcommand{\pt}[3]{\path (#1) node[croix]{} node[#2,font=\small,inner sep=2.5pt]{$#3$};}

\newcommand{\lycee@repere}[4]{%
  \draw[grille,step=1] (#1,#2) grid (#3,#4);
  \draw[-{Stealth[length=2mm]},thick] (#1,0)--({#3+0.4},0) node[below left=-1pt]{$x$};
  \draw[-{Stealth[length=2mm]},thick] (0,#2)--(0,{#4+0.4}) node[below left=-1pt]{$y$};
  \node[below left,font=\scriptsize,inner sep=1.5pt] at (0,0) {$0$};}
\newcommand{\lycee@gradx}[1]{\foreach \k/\l in {#1}{%
  \draw (\k,0.07)--(\k,-0.07) node[below,font=\scriptsize,inner sep=1.5pt]{$\l$};}}
\newcommand{\lycee@grady}[1]{\foreach \k/\l in {#1}{%
  \draw (0.07,\k)--(-0.07,\k) node[left,font=\scriptsize,inner sep=1.5pt]{$\l$};}}
\newcommand{\lycee@intff}[2]{\left[#1\,;#2\right]}
\newcommand{\lycee@intfo}[2]{\left[#1\,;#2\right[}
\newcommand{\lycee@intof}[2]{\left]#1\,;#2\right]}
\newcommand{\lycee@intoo}[2]{\left]#1\,;#2\right[}
\newcommand{\lycee@ens}[1]{\left\{#1\right\}}
\AtBeginDocument{%
  \providecommand{\repere}{\lycee@repere}%
  \providecommand{\gradx}{\lycee@gradx}%
  \providecommand{\grady}{\lycee@grady}%
  \providecommand{\Cf}{\mathcal{C}\sb{\scriptscriptstyle f}}%
  \providecommand{\Cg}{\mathcal{C}\sb{\scriptscriptstyle g}}%
  \providecommand{\intff}{\lycee@intff}%
  \providecommand{\intfo}{\lycee@intfo}%
  \providecommand{\intof}{\lycee@intof}%
  \providecommand{\intoo}{\lycee@intoo}%
  \providecommand{\ens}{\lycee@ens}}

\newlist{questions}{enumerate}{3}
\setlist[questions,1]{label=\arabic*\textdegree),leftmargin=*,itemsep=2pt,topsep=3pt}
\setlist[questions,2]{label=\alph*),leftmargin=*,itemsep=1pt,topsep=2pt}
\setlist[questions,3]{label=\roman*),leftmargin=*,itemsep=1pt,topsep=1pt}
\setlist[itemize]{label=\textbullet,leftmargin=*,itemsep=2pt,topsep=3pt}

\newcommand{\besoinplace}[1]{\par\penalty\@M\begingroup
  \dimen@=#1\relax\dimen@ii=\pagegoal\advance\dimen@ii-\pagetotal
  \ifdim\dimen@>\dimen@ii\ifdim\dimen@ii>\z@\vfil\fi\break\fi\endgroup}

\newcommand{\lycee@pied}{\thepage/\pageref{LastPage}}
\newcommand{\entete}[2]{%
  \pagestyle{fancy}\fancyhf{}%
  \fancyhead[L]{\footnotesize\color{coulPart}\textbf{#1}}%
  \fancyhead[R]{\footnotesize\color{coulPart}\textbf{#2}}%
  \fancyfoot[C]{\footnotesize\lycee@pied}%
  \renewcommand{\headrulewidth}{0.4pt}%
  \renewcommand{\headrule}{\hbox to\headwidth{%
    \color{coulPart!40}\leaders\hrule height \headrulewidth\hfill}}}

\RequirePackage{ccicons}
\newcommand{\lycee@licence}{{\scriptsize\color{gray}%
  \copyright\ Luc Destombe --- \ccbyncsa\ CC BY-NC-SA 4.0}}
\newif\iflycee@licence \lycee@licencetrue
\newcommand{\sanslicence}{\lycee@licencefalse}
\AtBeginDocument{\iflycee@licence\AddToHookNext{shipout/foreground}{%
  \put(0,0){\rlap{\kern\dimexpr1in+\hoffset+\oddsidemargin+\textwidth\relax
    \llap{\raisebox{-\dimexpr1in+\voffset+\topmargin+\headheight+\headsep
      +\textheight+\footskip\relax}{\lycee@licence}}}}}\fi}

\newcommand{\lycee@site}{faire-des-maths.fr}
\newcommand{\lycee@version}{{\scriptsize\color{gray}\lycee@site
  \ifdefined\lyceeversion\ --- Version~\lyceeversion\fi}}
\AtBeginDocument{\iflycee@licence\AddToHookNext{shipout/foreground}{%
  \put(0,0){\rlap{\kern\dimexpr1in+\hoffset+\oddsidemargin\relax
    \raisebox{-\dimexpr1in+\voffset+\topmargin+\headheight+\headsep
      +\textheight+\footskip\relax}{\lycee@version}}}}\fi}

\newcommand{\formulesserrees}{%
  \setlength{\abovedisplayskip}{3pt}\setlength{\belowdisplayskip}{3pt}%
  \setlength{\abovedisplayshortskip}{2pt}\setlength{\belowdisplayshortskip}{3pt}}

\setlength{\parindent}{0pt}

\RequirePackage[implicit=false,hidelinks,pdfencoding=auto,psdextra,bookmarksopen,
  bookmarksopenlevel=1,pdfcreator={},pdfproducer={}]{hyperref}
\RequirePackage{bookmark}
\hypersetup{pdfauthor={Luc Destombe},pdfkeywords={CC BY-NC-SA 4.0}}
\newcounter{lycee@signet}
\newcommand{\signet}[3][]{\stepcounter{lycee@signet}%
  \edef\lycee@dest{\if\relax\detokenize{#1}\relax
    signet.\arabic{lycee@signet}\else#1\fi}%
  \hypertarget{\lycee@dest}{}\bookmark[level=#2,dest=\lycee@dest]{#3}}
\pdfstringdefDisableCommands{%
  \def\R{R}\def\vect#1{#1}\def\Cf{Cf}\def\Cg{Cg}\def\fois#1{×#1}%
  \def\barre#1{#1}\def\intervalle#1#2{[#1;#2]}\def\intff#1#2{[#1;#2]}%
  \def\textsuperscript#1{#1}\def\dfrac#1#2{#1/#2}\def\frac#1#2{#1/#2}%
  \def\vec#1{#1⃗}}%
\RequirePackage[official]{eurosym}

\geometry{top=0.8cm,bottom=1.3cm,left=1cm,right=1cm,footskip=0.6cm}

\pagestyle{fancy}
\fancyhf{}
\renewcommand{\headrulewidth}{0pt}
\fancyfoot[C]{\footnotesize Page \thepage{} / \pageref{LastPage}}

\newcommand{\titrefiche}[2]{%
  \begin{center}
    \Large\textbf{#1}\\[2pt]
    \large\textbf{#2}\\[2pt]
  \end{center}
  \vspace{0.10cm}}

\setlength{\columnseprule}{0.4pt}
\setlength{\columnsep}{15pt}
\renewcommand{\columnseprulecolor}{\color{black!45}}
\setlength{\multicolsep}{2pt}

\newcounter{partiefiche}
\renewcommand{\thepartiefiche}{\Roman{partiefiche}}
\newif\ifapresbandeau
\newcommand{\lycee@bandeau}[1]{%
  \par\penalty-600
  \vspace{0.08cm}%
  \noindent\colorbox{coulPart}{%
    \parbox{\dimexpr\linewidth-2\fboxsep\relax}{%
      \color{white}\bfseries\raggedright\strut#1\strut}}%
  \nopagebreak\par\nopagebreak
  \vspace{0.06cm}\nopagebreak
  \apresbandeautrue}
\newcommand{\partiefiche}[1]{%
  \refstepcounter{partiefiche}%
  \lycee@bandeau{\signet{1}{\thepartiefiche{} --- #1}\thepartiefiche{} --- #1}}
\newcommand{\partiefichesansnum}[1]{\lycee@bandeau{\signet{1}{#1}#1}}

\newcounter{exo}
\newlength{\espaceexo}\setlength{\espaceexo}{7pt}
\newlength{\espacetitre}\setlength{\espacetitre}{2pt}
\newcommand{\exo}[1][\relax]{%
  \par
  \ifapresbandeau\apresbandeaufalse\else\penalty-150\addvspace{\espaceexo}\fi
  \refstepcounter{exo}%
  \noindent\signet[exercice-\arabic{exo}]{2}{Exercice \arabic{exo}}%
  \textbf{\color{coulExo}\underline{Exercice \theexo}}%
  \ifx\relax#1\relax\else\textbf{ : #1}\fi
  \par\nopagebreak\vspace{\espacetitre}\nopagebreak
  \ignorespaces}

\setlist[enumerate]{nosep,leftmargin=*,itemsep=2pt,topsep=2pt}
\setlist[itemize]{nosep,leftmargin=*,topsep=2pt}
\newcommand{\choix}[3]{\par\hspace*{1em}%
  \makebox[0.3\linewidth][l]{a) #1}\makebox[0.3\linewidth][l]{b) #2}c) #3}
\newcommand{\deux}[2]{\par\hspace*{0.6em}%
  \makebox[0.48\linewidth][l]{#1}#2}
\newcommand{\trois}[3]{\par\hspace*{0.6em}%
  \makebox[0.32\linewidth][l]{#1}\makebox[0.32\linewidth][l]{#2}#3}

\renewenvironment{center}{\par\addvspace{3pt}\centering}{\par\addvspace{3pt}}
\formulesserrees
\AtBeginDocument{\formulesserrees}
\clubpenalty=10000
\widowpenalty=10000
\displaywidowpenalty=10000
\brokenpenalty=10000
\setlength{\parskip}{0pt}

\RequirePackage{ragged2e}
\setlength{\RaggedRightRightskip}{0pt plus 1fil}
\AtBeginDocument{\RaggedRight\binoppenalty=10000 \relpenalty=10000 }
\g@addto@macro\@parboxrestore{\RaggedRight}
\makeatother
\geometry{bottom=1cm}

\setlength{\espaceexo}{4pt}
\setlist[enumerate]{nosep,leftmargin=*,itemsep=1pt,topsep=1pt}

\ShorthandsOff
\tikzset{
  vecteur/.style={-{Stealth[length=2.4mm]},line width=1pt,coulCourbe},
  vecteurB/.style={-{Stealth[length=2.4mm]},line width=1pt,coulCourbeB},
  vecteurC/.style={-{Stealth[length=2.4mm]},line width=1pt,coulCourbeC},
  base/.style={-{Stealth[length=2mm]},line width=0.9pt,black},
}
\ShorthandsOn
\newcommand{\coord}[2]{\begin{pmatrix}#1\\#2\end{pmatrix}}

\begin{document}

\titrefiche{Seconde --- Repérage, vecteurs et trigonométrie}{Fiche d'exercices du chapitre}

{\small\itshape Le repère $(O\,;\vec{\imath},\vec{\jmath})$ est orthonormé dans toute la
fiche ; on trace les repères sur le cahier, un carreau pour unité.\par}
\vspace{0.12cm}

\begin{multicols}{2}\raggedcolumns

\partiefiche{Repère du plan, base orthonormée}

\exo[Lire et placer]
\begin{center}
\begin{tikzpicture}[x=0.42cm,y=0.42cm]
  \repere{-5}{-5}{5}{4}
  \gradx{-4,-3,-2,-1,2,3,4} \grady{-4,-3,-2,-1,2,3}
  \draw[base] (0,0) -- node[midway,below,font=\footnotesize]{$\vec{\imath}$} (1,0);
  \draw[base] (0,0) -- node[midway,left,font=\footnotesize]{$\vec{\jmath}$} (0,1);
  \pt{2,3}{above right}{A} \pt{-3,2}{above left}{B} \pt{-4,-2}{below left}{C}
  \pt{1,-3}{below right}{D} \pt{0,2}{right}{E}
\end{tikzpicture}
\end{center}
\begin{enumerate}[label=\alph*)]
  \item Lire les coordonnées des points $A$, $B$, $C$, $D$ et $E$.
  \item Placer les points $F(3\,;-1)$, $G(-2\,;0)$ et $H(0\,;-4)$.
  \item Écrire $\vect{OA}$ et $\vect{OC}$ en fonction de $\vec{\imath}$ et $\vec{\jmath}$.
  \item Écrire $\vect{CA}$ en fonction de $\vec{\imath}$ et $\vec{\jmath}$.
\end{enumerate}

\exo[Abscisse et ordonnée]
Pour chaque point $K(-2\,;5)$, $L(0\,;3)$ et $M(4\,;0)$, écrire une phrase qui donne son
abscisse et son ordonnée, puis dire s'il est sur l'un des axes, et lequel.

\exo[Points définis par $\vec{\imath}$ et $\vec{\jmath}$]
Placer dans un repère les points $R$, $S$ et $T$ tels que
\begin{center}
$\vect{OR}=3\vec{\imath}-2\vec{\jmath}$ \qquad $\vect{OS}=-\vec{\imath}+4\vec{\jmath}$
\qquad $\vect{OT}=-2\vec{\jmath}$
\end{center}
puis donner leurs coordonnées.

\exo[Le repère d'un parallélogramme]
$ABCD$ est un parallélogramme de centre $K$. On se place dans le repère
$(A\,;\vect{AB},\vect{AD})$, qui n'est pas orthonormé.
\begin{enumerate}[label=\alph*)]
  \item Faire une figure.
  \item Donner les coordonnées de $A$, de $B$ et de $D$.
  \item Écrire $\vect{AC}$ en fonction de $\vect{AB}$ et $\vect{AD}$ ; en déduire les
        coordonnées de $C$.
  \item $K$ est le milieu de $[AC]$ : en déduire les coordonnées de $K$.
\end{enumerate}

\partiefiche{Coordonnées d'un vecteur}

\exo[Lire des coordonnées]
\begin{center}
\begin{tikzpicture}[x=0.4cm,y=0.4cm]
  \quadrillage{-0.4}{-0.4}{17.4}{4.4}
  \draw[base] (0,0) -- node[midway,below,font=\footnotesize]{$\vec{\imath}$} (1,0);
  \draw[base] (0,0) -- node[midway,left,font=\footnotesize]{$\vec{\jmath}$} (0,1);
  \draw[vecteur] (1,3) -- node[midway,above,font=\footnotesize]{$\vec{a}$} (5,2);
  \draw[vecteurB] (9,4) -- node[midway,above left,font=\footnotesize]{$\vec{b}$} (6,1);
  \draw[vecteurC] (10,0) -- node[midway,right,font=\footnotesize]{$\vec{c}$} (10,4);
  \draw[vecteur] (16,1) -- node[midway,below,font=\footnotesize]{$\vec{d}$} (12,2);
  \draw[vecteurB] (14,4) -- node[midway,above,font=\footnotesize]{$\vec{e}$} (17,4);
\end{tikzpicture}
\end{center}
\begin{enumerate}[label=\alph*)]
  \item Lire les coordonnées des vecteurs $\vec{a}$, $\vec{b}$, $\vec{c}$, $\vec{d}$ et
        $\vec{e}$.
  \item Deux de ces vecteurs sont opposés : lesquels ? Le justifier par leurs coordonnées.
\end{enumerate}

\exo[Tracer des représentants]
On considère le point $A(-1\,;1)$ et le point $B(4\,;-2)$.
\begin{enumerate}[label=\alph*)]
  \item Tracer les représentants d'origine $A$ des vecteurs $\vec{u}\coord{3}{2}$,
        $\vec{v}\coord{-2}{3}$, $\vec{w}\coord{0}{-3}$ et $\vec{t}\coord{-4}{0}$.
  \item Tracer un représentant de $\vec{s}\coord{3}{-1}$ d'extrémité $B$.
\end{enumerate}

\exo[Calculer des coordonnées]
On considère les points $A(2\,;5)$, $B(6\,;1)$, $C(-3\,;4)$, $D(1\,;-2)$ et $E(-4\,;-1)$.
Calculer les coordonnées des vecteurs $\vect{AB}$, $\vect{CD}$, $\vect{DE}$, $\vect{BC}$ et
$\vect{EA}$, puis écrire $\vect{AB}$ en fonction de $\vec{\imath}$ et $\vec{\jmath}$.

\exo[Lire, puis calculer]
\begin{center}
\begin{tikzpicture}[x=0.42cm,y=0.42cm]
  \repere{-4}{-4}{4}{4}
  \gradx{-3,-2,-1,1,2,3} \grady{-3,-2,-1,1,2,3}
  \pt{-3,1}{above left}{A} \pt{1,3}{above right}{B}
  \pt{3,-1}{below right}{C} \pt{-1,-3}{below left}{D}
\end{tikzpicture}
\end{center}
\begin{enumerate}[label=\alph*)]
  \item Lire les coordonnées des points $A$, $B$, $C$ et $D$.
  \item Calculer les coordonnées des vecteurs $\vect{AB}$, $\vect{DC}$ et $\vect{BC}$.
  \item Vérifier ces résultats en comptant les carreaux sur la figure.
  \item Que remarque-t-on pour $\vect{AB}$ et $\vect{DC}$ ?
\end{enumerate}

\exo[Un point inconnu]
\begin{enumerate}[label=\alph*)]
  \item On considère le point $A(3\,;2)$. Calculer les coordonnées du point $E$ tel que
        $\vect{AE}\coord{-1}{4}$.
  \item On considère le point $B(-5\,;9)$. Calculer les coordonnées du point $F$ tel que
        $\vect{FB}\coord{2}{8}$.
\end{enumerate}

\partiefiche{Coordonnées d'une somme, d'un produit par un nombre réel}

\exo[Calculs de coordonnées]
On considère les vecteurs $\vec{u}\coord{2}{-3}$, $\vec{v}\coord{-4}{5}$ et
$\vec{w}\coord{3}{2}$. Calculer les coordonnées de :
\choix{$\vec{u}+\vec{v}$}{$3\vec{w}$}{$\vec{u}-2\vec{v}$}
\par\hspace*{1em}d)~$-0{,}5\vec{v}$ \qquad e)~$2\vec{u}+\vec{v}-\vec{w}$

\exo[Avec des points]
On considère les points $A(1\,;2)$, $B(4\,;-1)$ et $C(-2\,;3)$.
\begin{enumerate}[label=\alph*)]
  \item Calculer les coordonnées de $\vect{AB}$ et de $\vect{AC}$.
  \item Calculer les coordonnées de $\vect{AB}+\vect{AC}$ et de $2\vect{AB}-\vect{AC}$.
  \item Calculer les coordonnées de $\vect{BC}$, puis vérifier que
        $\vect{BC}=\vect{AC}-\vect{AB}$.
\end{enumerate}

\exo[Un point défini par une somme]
On considère le point $A(-1\,;3)$ et les vecteurs $\vec{u}\coord{2}{-3}$ et
$\vec{v}\coord{-1}{5}$.
\begin{enumerate}[label=\alph*)]
  \item Calculer les coordonnées de $\vec{u}+\vec{v}$ et de $2\vec{u}-\vec{v}$.
  \item En déduire les coordonnées des points $E$ et $F$ tels que
        $\vect{AE}=\vec{u}+\vec{v}$ et $\vect{AF}=2\vec{u}-\vec{v}$.
  \item Calculer les coordonnées du vecteur $\vect{EF}$.
\end{enumerate}

\exo[Démontrer une égalité]
On considère les vecteurs $\vec{u}\coord{2}{1}$, $\vec{v}\coord{1}{-2}$ et
$\vec{w}\coord{3}{4}$.
\begin{enumerate}[label=\alph*)]
  \item Démontrer que $2\vec{u}=\vec{v}+\vec{w}$.
  \item Sur un quadrillage, tracer depuis un même point un représentant de $2\vec{u}$,
        puis un représentant de $\vec{v}$ suivi d'un représentant de $\vec{w}$ : que
        constate-t-on ?
\end{enumerate}

\partiefiche{Combinaison de deux vecteurs}

\exo[Lire une combinaison]
\begin{center}
\begin{tikzpicture}[x=0.38cm,y=0.38cm]
  \quadrillage{-0.4}{-0.4}{18.4}{5.4}
  \draw[vecteur] (0,0) -- node[midway,below right,font=\footnotesize]{$\vec{a}$} (3,1);
  \draw[vecteurB] (1,2) -- node[midway,left,font=\footnotesize]{$\vec{b}$} (2,5);
  \draw[vecteurC] (4,0) -- node[midway,left,font=\footnotesize]{$\vec{w}_{1}$} (8,4);
  \draw[vecteurC] (9,3) -- node[midway,above,font=\footnotesize]{$\vec{w}_{2}$} (14,2);
  \draw[vecteurC] (17,0) -- node[midway,right,font=\footnotesize]{$\vec{w}_{3}$} (16,5);
\end{tikzpicture}
\end{center}
Écrire $\vec{w}_{1}$, $\vec{w}_{2}$ et $\vec{w}_{3}$ en fonction de $\vec{a}$ et $\vec{b}$.

\exo[Construire une combinaison]
\begin{center}
\begin{tikzpicture}[x=0.38cm,y=0.38cm]
  \quadrillage{-0.4}{-0.4}{14.4}{6.4}
  \draw[vecteur] (0,4) -- node[midway,above,font=\footnotesize]{$\vec{c}$} (4,5);
  \draw[vecteurB] (1,0) -- node[midway,left,font=\footnotesize]{$\vec{d}$} (0,2);
  \pt{3,0}{below}{M}
\end{tikzpicture}
\end{center}
On pose $\vec{r}=\vec{c}+2\vec{d}$ et $\vec{s}=2\vec{c}-\vec{d}$. Reproduire la figure,
puis construire un représentant de $\vec{r}$ et un représentant de $\vec{s}$, tous deux
d'origine $M$.

\exo[Dans un parallélogramme, dans un triangle]
\begin{enumerate}[label=\alph*)]
  \item $ABCD$ est un parallélogramme de centre $K$. Écrire $\vect{BD}$, puis $\vect{BK}$,
        en fonction de $\vect{AB}$ et $\vect{AD}$.
  \item $EFG$ est un triangle et $M$ le point de $[FG]$ tel que
        $\vect{FM}=\frac{2}{3}\vect{FG}$. Écrire $\vect{EM}$ en fonction de $\vect{EF}$ et
        $\vect{EG}$.
  \item En déduire les coordonnées de $M$ dans le repère $(E\,;\vect{EF},\vect{EG})$.
\end{enumerate}

\exo[Combinaisons et coordonnées]
On considère les vecteurs $\vec{a}\coord{2}{1}$ et $\vec{b}\coord{-1}{3}$.
\begin{enumerate}[label=\alph*)]
  \item Calculer les coordonnées de $2\vec{a}+\vec{b}$ et de $\vec{a}-\vec{b}$.
  \item On considère le vecteur $\vec{w}\coord{0}{7}$. Vérifier que
        $\vec{w}=\vec{a}+2\vec{b}$.
  \item Sur un quadrillage, tracer $\vec{w}$, puis le chemin formé de $\vec{a}$ et de
        deux fois $\vec{b}$ qui va de son origine à sa pointe.
  \item On considère le point $A(1\,;1)$. Calculer les coordonnées du point $M$ tel que
        $\vect{AM}=\vec{a}-2\vec{b}$.
\end{enumerate}

\partiefiche{Égalité de vecteurs et parallélogrammes}

\exo[Parallélogramme ou non]
Pour chaque quadrilatère, dire s'il est un parallélogramme.
\begin{enumerate}[label=\alph*)]
  \item $ABCD$, avec $A(-2\,;1)$, $B(1\,;3)$, $C(4\,;1)$ et $D(1\,;-1)$.
  \item $KLMN$, avec $K(-1\,;-2)$, $L(2\,;-1)$, $M(4\,;3)$ et $N(1\,;2)$.
  \item $RSTU$, avec $R(0\,;1)$, $S(3\,;2)$, $T(4\,;5)$ et $U(0\,;4)$.
\end{enumerate}

\exo[Le quatrième sommet]
On considère les points $A(-2\,;1)$, $B(0\,;4)$ et $C(3\,;2)$.
\begin{enumerate}[label=\alph*)]
  \item Calculer les coordonnées du point $D$ tel que $ABCD$ soit un parallélogramme.
  \item Calculer les coordonnées du point $E$ tel que $ABEC$ soit un parallélogramme.
  \item Placer les cinq points et tracer les deux parallélogrammes.
\end{enumerate}

\exo[Image par une translation]
On considère les points $A(1\,;-1)$, $B(2\,;-3)$ et $C(4\,;2)$. Le point $F$ est l'image de
$B$ par la translation de vecteur $\vect{AC}$.
\begin{enumerate}[label=\alph*)]
  \item Écrire l'égalité de vecteurs qui définit $F$, puis calculer ses coordonnées.
  \item Démontrer que $ABFC$ est un parallélogramme.
  \item Écrire une phrase qui dit de quel point $C$ est l'image, et par quelle
        translation.
  \item Calculer les coordonnées du point $G$ tel que $\vect{AG}=\vect{AB}+\vect{AC}$ :
        que remarque-t-on ?
\end{enumerate}

\partiefiche{Colinéarité, déterminant}

\exo[Colinéaires ou non]
Pour chaque couple de vecteurs, dire si les vecteurs sont colinéaires.
\deux{a)~$\vec{u}\coord{4}{-6}$ et $\vec{v}\coord{-6}{9}$}{b)~$\vec{u}\coord{2}{5}$ et
$\vec{v}\coord{5}{2}$}
\deux{c)~$\vec{u}\coord{\frac{2}{3}}{1}$ et $\vec{v}\coord{2}{3}$}{d)~$\vec{u}\coord{3}{-1}$
et $\vec{v}\coord{-5}{2}$}

\exo[Trouver une coordonnée]
On considère le vecteur $\vec{u}\coord{-2}{4}$.
\begin{enumerate}[label=\alph*)]
  \item Donner les coordonnées de deux vecteurs colinéaires à $\vec{u}$.
  \item Calculer le nombre réel $x$ tel que $\vec{v}\coord{3}{x}$ soit colinéaire à
        $\vec{u}$.
\end{enumerate}

\exo[Parallèles ou non]
\begin{enumerate}[label=\alph*)]
  \item On considère les points $A(1\,;2)$, $B(4\,;3)$, $C(-2\,;-1)$ et $D(4\,;1)$. Les
        droites $(AB)$ et $(CD)$ sont-elles parallèles ? Le quadrilatère $ABDC$
        est-il un parallélogramme ?
  \item Même première question pour les droites $(EF)$ et $(GH)$, avec $E(0\,;3)$, $F(2\,;-1)$,
        $G(-1\,;0)$ et $H(2\,;-5)$.
\end{enumerate}

\exo[Alignés ou non]
Pour chaque cas, dire si les trois points sont alignés.
\begin{enumerate}[label=\alph*)]
  \item $A(1\,;3)$, $B(-1\,;-1)$ et $C(4\,;9)$.
  \item $D(3\,;7)$, $E(1\,;2)$ et $F(7\,;16)$.
  \item Le point $G(10\,;-6)$ est-il sur la droite $(KL)$, avec $K(4\,;-2)$ et
        $L(-2\,;2)$ ?
  \item Avec les points du a), calculer les coordonnées du point $M$ tel que
        $\vect{AM}=\vect{AB}+\vect{AC}$. Le point $M$ est-il sur la droite $(AB)$ ?
\end{enumerate}

\partiefiche{Milieu d'un segment}

\exo[Calculer des milieux]
Pour chaque segment $[AB]$, calculer les coordonnées de son milieu.
\deux{a)~$A(2\,;5)$ et $B(6\,;-1)$}{b)~$A(10\,;1)$ et $B(-4\,;5)$}
\deux{c)~$A(3\,;-4)$ et $B(-3\,;4)$}{d)~$A(-1\,;2)$ et $B(4\,;7)$}

\exo[Retrouver une extrémité]
Le point $I(1\,;1)$ est le milieu de $[AB]$. Calculer les coordonnées de $A$ quand :
\deux{a)~$B(5\,;-3)$}{b)~$B(-2\,;4)$}

\exo[De deux façons]
On considère les points $E(-2\,;5)$, $F(1\,;2)$, $G(4\,;-1)$ et $H(7\,;-4)$.
\begin{enumerate}[label=\alph*)]
  \item Calculer $\vect{EF}$ et $\vect{FG}$. Que peut-on en déduire pour le point $F$ ?
  \item Retrouver ce résultat avec les coordonnées du milieu de $[EG]$.
  \item Le point $G$ est-il le milieu de $[FH]$ ?
  \item Démontrer, avec un déterminant, que les points $E$, $F$ et $H$ sont alignés.
\end{enumerate}

\exo[La droite des milieux]
On considère les points $A(-2\,;1)$, $B(4\,;3)$ et $C(2\,;-3)$ ; $I$ est le milieu de
$[AB]$ et $J$ celui de $[AC]$.
\begin{enumerate}[label=\alph*)]
  \item Calculer les coordonnées de $I$ et de $J$.
  \item Démontrer que les droites $(IJ)$ et $(BC)$ sont parallèles.
  \item Comparer les coordonnées de $\vect{IJ}$ et de $\vect{BC}$ ; que remarque-t-on ?
\end{enumerate}

\partiefiche{Norme d'un vecteur, longueur d'un segment}

\exo[Des normes]
Calculer la norme de chaque vecteur $\vec{u}\coord{0}{-1}$, $\vec{v}\coord{3}{4}$,
$\vec{w}\coord{8}{-6}$, $\vec{t}\coord{2}{-2}$ et $\vec{s}\coord{-1}{3}$ (valeurs exactes,
simplifiées).

\exo[Des longueurs]
Calculer les longueurs $AB$, $CD$ et $EF$ (valeurs exactes, simplifiées), avec
$A(1\,;2)$, $B(4\,;6)$, $C(-2\,;3)$, $D(4\,;-5)$, $E(-1\,;-1)$ et $F(3\,;1)$.

\exo[Nature d'un triangle]
\begin{enumerate}[label=\alph*)]
  \item On considère les points $A(-1\,;2)$, $B(3\,;5)$ et $C(6\,;1)$. Démontrer que le
        triangle $ABC$ est rectangle isocèle.
  \item Calculer les coordonnées du point $D$ tel que $ABCD$ soit un parallélogramme.
        Quelle est la nature de $ABCD$ ?
  \item On considère les points $E(1\,;1)$, $F(5\,;2)$ et $G(2\,;5)$. Le triangle $EFG$
        est-il isocèle ? rectangle ?
\end{enumerate}

\exo[Sur un cercle]
$\mathcal{C}$ est le cercle de centre $K(1\,;-1)$ et de rayon $5$. Pour chaque point
$M(4\,;3)$, $N(-2\,;-5)$, $R(6\,;-1)$ et $S(0\,;4)$, dire s'il est sur le cercle
$\mathcal{C}$. Calculer ensuite les coordonnées du milieu de $[MN]$ : que dire du segment
$[MN]$ ?

\partiefiche{Trigonométrie dans le triangle rectangle}

\exo[Calculer une longueur]
Donner les longueurs arrondies au dixième.
\begin{enumerate}[label=\alph*)]
  \item $ABC$ est rectangle en $B$, avec $\widehat{A}=40^{\circ}$ et $AC=10$~cm.
        Calculer $BC$, puis $AB$.
  \item $DEF$ est rectangle en $E$, avec $\widehat{D}=25^{\circ}$ et $DE=6$~cm.
        Calculer $EF$, puis $DF$.
\end{enumerate}

\exo[Les côtés d'un triangle rectangle]
$RST$ est un triangle rectangle en $S$.
\begin{enumerate}[label=\alph*)]
  \item Écrire une phrase qui nomme l'hypoténuse, puis le côté adjacent et le côté
        opposé à l'angle $\widehat{R}$.
  \item Même question pour l'angle $\widehat{T}$.
  \item Écrire $\cos\widehat{R}$, $\sin\widehat{R}$ et $\tan\widehat{T}$ comme des
        quotients de longueurs.
\end{enumerate}

\exo[Calculer un angle]
Donner les angles arrondis au degré.
\begin{enumerate}[label=\alph*)]
  \item $GHK$ est rectangle en $H$, avec $GH=5$~cm et $GK=8$~cm. Calculer $\widehat{G}$.
  \item $LMN$ est rectangle en $M$, avec $LM=4$~cm et $MN=7$~cm. Calculer $\widehat{L}$.
  \item Dans un repère, on considère $R(1\,;1)$, $S(5\,;1)$ et $T(5\,;4)$ ; le triangle
        $RST$ est rectangle en $S$. Calculer la longueur $RT$, puis l'angle
        $\widehat{R}$.
\end{enumerate}

\exo[Une échelle]
Une échelle de $5$~m est posée contre un mur vertical, sur un sol horizontal. Pour
travailler en sécurité, l'angle entre l'échelle et le sol doit être compris entre
$65^{\circ}$ et $75^{\circ}$.
\begin{enumerate}[label=\alph*)]
  \item L'échelle fait un angle de $70^{\circ}$ avec le sol. À quelle hauteur arrive-t-elle
        sur le mur ? À quelle distance du mur est son pied ? (au dixième)
  \item Le pied de l'échelle est à $1{,}2$~m du mur. Calculer l'angle avec le sol, au
        degré : l'échelle est-elle posée en sécurité ?
\end{enumerate}

\exo[Une rampe d'accès]
Une rampe d'accès pour fauteuils roulants doit monter de $0{,}6$~m jusqu'à l'entrée
d'une mairie. L'angle entre la rampe et le sol ne doit pas dépasser $3^{\circ}$.
\begin{enumerate}[label=\alph*)]
  \item Une rampe de $8$~m de long à l'horizontale convient-elle ? (calculer son angle)
  \item Calculer la longueur horizontale minimale de la rampe, au dixième de mètre.
\end{enumerate}

\end{multicols}

\end{document}
